%% kapa: exam=1 round=32 year=2021 period=2 form=A subject=appraisal_law subject_ko=감정평가관계법규 seq=3 number=3
%% source: Q-Net 감정평가사 기출문제 (cst00309) artlSeq=5208397 file=extracted/1008/5208397/A형/2021년도 제32회 감정평가사 1차 2교시 A형.pdf
%% conversion: kordoc 4.24.51 -> tools/gen1.py (auto); review=proofread-sonnet; flags=none
\begin{question}{3}
국토의 계획 및 이용에 관한 법령상 도시·군기본계획에 관한 설명으로 옳은 것은?
\begin{choices}
  \item 특별시장·광역시장·특별자치시장·도지사·특별자치도지사는 관할 구역에 대하여 도시·군기본계획을 수립하여야 한다.
  \item 시장 또는 군수가 도시·군기본계획을 변경하려면 지방의회의 승인을 받아야 한다.
  \item 도시·군기본계획을 변경하기 위하여 공청회를 개최한 경우, 공청회에서 제시된 의견이 타당하다고 인정하더라도 도시·군기본계획에 반영하지 않을 수 있다.
  \item 도시·군기본계획 입안일부터 5년 이내에 토지적성평가를 실시한 경우에는 도시·군기본계획의 수립을 위한 기초조사의 내용에 포함되어야 하는 토지적성평가를 하지 아니할 수 있다.
  \item 도지사는 시장 또는 군수가 수립한 도시·군기본계획에 대하여 관계 행정기관의 장과 협의하였다면, 지방도시계획위원회의 심의를 거치지 아니하고 승인할 수 있다.
\end{choices}
\end{question}
\answer{4}
